\section*{Capítulo \ref{ch:g8:ice-on-earth} --- O gelo na Terra}

\begin{solution}{exo:g8:ice-on-earth:1}
(a) A zona de acumulação. (b) A zona de ablação. (c) A extremidade.
(d) Uma \omterm{def:g8:ice-on-earth:glacial-landforms}{morena} (frontal).
\end{solution}

\begin{solution}{exo:g8:ice-on-earth:2}
(a) Falso: uma \omterm{def:g8:ice-on-earth:glacier}{geleira} é feita de \omterm{def:g4:clouds-rain-snow:precipitation}{neve} caída sobre a terra e transformada
em gelo; a água do \omterm{def:g2:water-on-earth:ocean}{mar} congelada é o \omterm{def:g8:ice-on-earth:sea-ice}{gelo marinho}. (b) Verdadeiro.
(c) Falso: o seu gelo continua descendo a encosta; a extremidade recua
porque o gelo derrete mais depressa do que desce. (d) Verdadeiro.
\end{solution}

\begin{solution}{exo:g8:ice-on-earth:3}
(a) Um \omterm{def:g8:ice-on-earth:glacial-landforms}{vale em U}. (b) Um bloco errático. (c) Uma \omterm{def:g8:ice-on-earth:ice-sheet}{plataforma de gelo}. (d) O \omterm{def:g8:ice-on-earth:permafrost}{permafrost}.
\end{solution}

\begin{solution}{exo:g8:ice-on-earth:4}
(a) A Antártida e a Groenlândia. (b) A Antártida, cerca de \qty{58}{m}
(Groenlândia: cerca de \qty{7.4}{m}). (c) Cerca de um décimo (10~\%).
\end{solution}

\begin{solution}{exo:g8:ice-on-earth:5}
(a) Cerca de \num{7.05}~milhões~de~\unit{km^2} em 1979 e \num{3.57}~milhões~de~\unit{km^2}
em 2012. (b) $7.05 - 3.57 = 3.48$: cerca de \num{3.5}~milhões~de~\unit{km^2}.
(c) $3.48 \div 7.05 \approx 0.49$: cerca da metade. (d) Um único ano
depende do \omterm{def:g1:watching-the-weather:weather}{tempo} daquele verão: 2012 fica bem abaixo da linha. A linha de
tendência usa todos os anos, então um verão fora do comum a muda muito
pouco.
\end{solution}

\begin{solution}{exo:g8:ice-on-earth:6}
(a) $2004 - 1941 = 63$ anos. (b) $12000 \div 63 \approx 190.5$: cerca de
\qty{190}{m} por ano. (c) $2021 - 1984 = 37$ anos; $3200 \div 37 \approx
86$: cerca de \qty{86}{m} por ano. (d) A \omterm{def:g8:ice-on-earth:glacier}{geleira} Muir, mais de duas vezes
mais depressa.
\end{solution}

\begin{solution}{exo:g8:ice-on-earth:7}
Na \emph{zona de acumulação}, no alto da \omterm{def:g1:landscapes-around-us:landform}{montanha}, cai mais \omterm{def:g4:clouds-rain-snow:precipitation}{neve} a cada
ano do que derrete. A \omterm{def:g4:clouds-rain-snow:precipitation}{neve} nova enterra e comprime a \omterm{def:g4:clouds-rain-snow:precipitation}{neve} velha, que vira
\emph{firn} depois de cerca de um ano, e depois, mais fundo e muito mais
tarde, gelo sólido. Sob o seu próprio peso, o gelo \emph{desce} devagar
pelo \omterm{def:g1:landscapes-around-us:landform}{vale}. Mais abaixo, na \emph{zona de ablação}, ele derrete mais
depressa do que cai a \omterm{def:g4:clouds-rain-snow:precipitation}{neve}; o gelo que chega à extremidade é um gelo que
se formou lá no alto, muitas vezes há muito tempo.
\end{solution}

\begin{solution}{exo:g8:ice-on-earth:8}
(a) No copo B: o gelo derretido escorre para a água e a faz subir; no copo
A, o nível fica onde estava. (b) O copo A é o \omterm{def:g8:ice-on-earth:ice-sheet}{iceberg} (gelo que já flutua
no \omterm{def:g2:water-on-earth:ocean}{mar}); o copo B é o \omterm{def:g8:ice-on-earth:ice-sheet}{manto de gelo} (gelo sobre a terra). (c) Uma
\omterm{def:g8:ice-on-earth:ice-sheet}{plataforma de gelo} segura as \omterm{def:g8:ice-on-earth:glacier}{geleiras} atrás dela. Quando ela se quebra,
essas \omterm{def:g8:ice-on-earth:glacier}{geleiras} aceleram e despejam mais gelo de terra no \omterm{def:g2:water-on-earth:ocean}{mar}, e isso o
faz subir.
\end{solution}

\begin{solution}{exo:g8:ice-on-earth:9}
(a) Um \omterm{def:g8:ice-on-earth:glacial-landforms}{vale em U}, escavado por uma \omterm{def:g8:ice-on-earth:glacier}{geleira}. (b) Uma \omterm{def:g8:ice-on-earth:glacier}{geleira} o trouxe lá
de cima no \omterm{def:g1:landscapes-around-us:landform}{vale}, onde há granito, e o deixou quando o gelo derreteu: é um
bloco errático. (c) Uma \omterm{def:g8:ice-on-earth:glacial-landforms}{morena} lateral: a sua altura mostra que espessura
a \omterm{def:g8:ice-on-earth:glacier}{geleira} teve um dia; o gelo chegou pelo menos até essa altura na
encosta do \omterm{def:g1:landscapes-around-us:landform}{vale}.
\end{solution}

\begin{solution}{exo:g8:ice-on-earth:10}
(a) $7.4 \div 10 = \qty{0.74}{m}$. (b) $58 \div 100 = \qty{0.58}{m}$.
(c) $58 \div 7.4 \approx 7.8$: cerca de 8 vezes mais.
\end{solution}

\begin{solution}{exo:g8:ice-on-earth:11}
(a) $7.21 \times 10^6 = \qty{7210000}{km^2}$ e $4.65 \times 10^6 =
\qty{4650000}{km^2}$. (b) $7.21 - 4.65 = \num{2.56}\text{ milhões de }\unit{km^2}$, ou
seja, \qty{2560000}{km^2}. (c) $2.56 \div 3 \approx 0.85$: cerca de
\num{0.85}~milhão~de~\unit{km^2} por década, um pouco mais que os
\num{0.76}~milhão~de~\unit{km^2} da linha de tendência; cada média de década ainda leva os altos e baixos
dos seus próprios dez verões, mas as duas maneiras de ler o gráfico mostram
a mesma grande perda. (d) A extensão não pode cair abaixo de zero, e uma
linha de tendência só descreve os anos que foram medidos: o que acontece
depois depende de como o \omterm{def:g5:weather-and-climate:climate}{clima} do Ártico muda, então a linha não é uma
previsão.
\end{solution}

\begin{solution}{exo:g8:ice-on-earth:12}
(a) $120 \div 1 = 120$: o \omterm{def:g1:landscapes-around-us:water}{litoral} ficava cerca de \qty{120}{km} mais longe,
\omterm{def:g2:water-on-earth:ocean}{mar} adentro. (b) Ela tinha evaporado do \omterm{def:g2:water-on-earth:ocean}{oceano}, caído como \omterm{def:g4:clouds-rain-snow:precipitation}{neve} sobre os
\omterm{def:g3:continents-and-oceans:continent}{continentes} e ficado ali como o gelo dos grandes \omterm{def:g8:ice-on-earth:ice-sheet}{mantos de gelo}. (c) Uma
\omterm{def:g8:ice-on-earth:glacial-landforms}{morena} frontal se forma na extremidade de uma \omterm{def:g8:ice-on-earth:glacier}{geleira}. Uma \omterm{def:g8:ice-on-earth:glacial-landforms}{morena} numa
\omterm{def:g1:landscapes-around-us:landform}{planície} longe de qualquer \omterm{def:g1:landscapes-around-us:landform}{montanha} mostra que um \omterm{def:g8:ice-on-earth:ice-sheet}{manto de gelo} chegou um
dia até essa \omterm{def:g1:landscapes-around-us:landform}{planície}, numa época em que o gelo cobria muito mais terra
que hoje.
\end{solution}

\begin{solution}{pb:g8:ice-on-earth:1}
\textbf{1.} $1080 - 0 = \qty{1080}{m}$.

\textbf{2.} 1960--1970: \qty{180}{m}; 1970--1980: \qty{150}{m};
1980--1990: \qty{190}{m}; 1990--2000: \qty{180}{m}; 2000--2010:
\qty{180}{m}; 2010--2020: \qty{200}{m}. Menos em 1970--1980, mais em
2010--2020.

\textbf{3.} Não: uma \omterm{def:g8:ice-on-earth:glacier}{geleira} que recua continua descendo a encosta. A sua
extremidade recua porque o gelo derrete mais depressa do que chega.

\textbf{4.} Uma \omterm{def:g8:ice-on-earth:glacial-landforms}{morena} frontal: o gelo levou detritos rochosos \omterm{def:g1:landscapes-around-us:landform}{vale} abaixo
e os largou na sua extremidade, onde o gelo derretia, enquanto a
extremidade ficava ali.

\textbf{5.} $1080 \div 6 = \qty{180}{m}$ por década.

\textbf{6.} $180 \div 10 = \qty{18}{m}$ por ano.

\textbf{7.} $190 \div 18 \approx 10.6$: cerca de 11 vezes mais rápido.

\textbf{8.} Menos: ela perdeu mais gelo na sua zona de ablação do que
ganhou na sua zona de acumulação, então a sua extremidade recuou.

\textbf{9.} Em U: um fundo largo com encostas íngremes, porque a \omterm{def:g8:ice-on-earth:glacier}{geleira}
alargou e aprofundou o \omterm{def:g1:landscapes-around-us:landform}{vale} enquanto o enchia.

\textbf{10.} $1980 - 1080 = \qty{900}{m}$.

\textbf{11.} $900 \div 18 = 50$ anos: em \textbf{2070}.

\textbf{12.} A velocidade mudou de uma década para outra (de 150 a
\qty{200}{m}), e o recuo futuro depende da queda de \omterm{def:g4:clouds-rain-snow:precipitation}{neve} e das
temperaturas de verão dos próximos anos: um ritmo passado não é uma
previsão.
\end{solution}
